% fixture: algorithm/here-after-heading
% tier: T3
% engine: document
% covers: [H] algorithms right after a heading: \@setnobreak clears \if@nobreak in the box and \@endfloatbox's \outer@nobreak sets it again, so the text after the box is still right after the heading (no indentation, \clubpenalty\@M); two in a row; after a run-in heading
% issue: #277
\documentclass{article}
\usepackage{algorithm}
\usepackage{algpseudocode}
\begin{document}
\section{Greatest common divisor}
\begin{algorithm}[H]
\caption{Euclid's algorithm, with a caption long enough to run over two lines of the text width}
\begin{algorithmic}[1]
\State $x \gets a$ and a statement long enough to run over two lines of the algorithm, as it goes on and on
\State $y \gets b$
\end{algorithmic}
\end{algorithm}
Each step at least halves one of the numbers every two steps, so the loop runs a logarithmic number of times, and more text.

Another paragraph that runs over more than one line of text, with words and more words in it, to see.
\subsection{Two in a row}
\begin{algorithm}[H]
\begin{algorithmic}
\State $x \gets a$
\end{algorithmic}
\end{algorithm}
\begin{algorithm}[H]
\begin{algorithmic}
\State $y \gets b$
\end{algorithmic}
\end{algorithm}
Text after them, long enough to run over two lines of the page's text width, with more words in it here.

\paragraph{Run-in.}
\begin{algorithm}[H]
\begin{algorithmic}
\State $x \gets a$
\end{algorithmic}
\end{algorithm}
Text after it, long enough to run over two lines of the page's text width, with more words in it here.
\end{document}
